Activation steering and prompt engineering are widely used to control language model behavior, but they rely on the target behavior being expressible in natural language. Many practically important features---such as making text sound less AI-generated, controlling hedging patterns, or managing sycophancy---resist precise verbalization. We hypothesize that learned token embeddings (\emph{concept tokens}) provide the greatest advantage over explicit prompting for features that are hardest to describe. We test this on \pythia across six behavioral features spanning a range of describability, from easy (sentiment, formality) to hard (AI-sounding text, hedging). For each feature, we extract a residual stream direction via contrastive mean difference, measure how consistently natural language prompts activate it, and train a concept token embedding to reference it. We find that concept tokens consistently outperform prompting for hard-to-describe features: AI-sounding text (+106\%), sycophancy (+60\%), and hedging (+25\%), while prompting performs adequately for easy features. The Spearman correlation between feature difficulty and concept token advantage is $\rho = 0.71$, a large effect size. These results suggest that learned embeddings can bridge the gap between features that are linearly represented in the residual stream and those that can be controlled via natural language, opening new possibilities for manipulating model behaviors that resist prompt engineering.
